%%%%%%%%%%%%%%%%%%%%%%%%%%%%%%%%%%%%%%%%%
% Short Sectioned Assignment LaTeX Template Version 1.0 (5/5/12)
% This template has been downloaded from: http://www.LaTeXTemplates.com
% Original author:  Frits Wenneker (http://www.howtotex.com)
% License: CC BY-NC-SA 3.0 (http://creativecommons.org/licenses/by-nc-sa/3.0/)
%%%%%%%%%%%%%%%%%%%%%%%%%%%%%%%%%%%%%%%%%

% \documentclass[paper=a4, fontsize=11pt]{scrartcl} % A4 paper and 11pt font size
\documentclass[11pt, a4paper]{report}
\usepackage[T1]{fontenc} % Use 8-bit encoding that has 256 glyphs
\usepackage[utf8]{inputenc}
\usepackage{fourier} % Use the Adobe Utopia font for the document - comment this line to return to the LaTeX default
\usepackage{listings} % para insertar código con formato similar al editor
\usepackage[spanish, es-tabla]{babel} % Selecciona el español para palabras introducidas automáticamente, p.ej. "septiembre" en la fecha y especifica que se use la palabra Tabla en vez de Cuadro
\usepackage{url} % ,href} %para incluir URLs e hipervínculos dentro del texto (aunque hay que instalar href)
\usepackage{graphics,graphicx, float} %para incluir imágenes y colocarlas
\usepackage[gen]{eurosym} %para incluir el símbolo del euro
\usepackage{cite} %para incluir citas del archivo <nombre>.bib
\usepackage{enumerate}
\usepackage{hyperref}
\hypersetup{
	colorlinks=true,	% false: boxed links; true: colored links
	linkcolor=black,	% color of internal links
	urlcolor=cyan		% color of external links
}
\renewcommand{\familydefault}{\sfdefault}
\usepackage{fancyhdr} % Custom headers and footers
\pagestyle{fancyplain} % Makes all pages in the document conform to the custom headers and footers
\fancyhead[L]{} % Empty left header
\fancyhead[C]{} % Empty center header
\fancyhead[R]{Sergio Cervilla Ortega} % My name
\fancyfoot[L]{} % Empty left footer
\fancyfoot[C]{} % Empty center footer
\fancyfoot[R]{\thepage} % Page numbering for right footer
%\renewcommand{\headrulewidth}{0pt} % Remove header underlines
\renewcommand{\footrulewidth}{0pt} % Remove footer underlines
\setlength{\headheight}{13.6pt} % Customize the height of the header

\begin{document}

	% Plantilla portada UGR
	\begin{titlepage}
\newlength{\centeroffset}
\setlength{\centeroffset}{-0.5\oddsidemargin}
\addtolength{\centeroffset}{0.5\evensidemargin}
\thispagestyle{empty}

\noindent\hspace*{\centeroffset}\begin{minipage}{\textwidth}

\centering
\includegraphics[width=0.9\textwidth]{imagenes/logo_ugr.jpg}\\[1.4cm]

\textsc{ \Large TRABAJO FIN DE GRADO\\[0.2cm]}
\textsc{ GRADO EN INGENIERIA INFORMATICA}\\[1cm]

{\Huge\bfseries Chief \\}
\noindent\rule[-1ex]{\textwidth}{3pt}\\[3.5ex]
{\large\bfseries Apoyo a la labor de la Policía Local }
\end{minipage}

\vspace{2.5cm}
\noindent\hspace*{\centeroffset}
\begin{minipage}{\textwidth}
\centering

\textbf{Autor}\\ {Sergio Cervilla Ortega}\\[2.5ex]
\textbf{Director}\\ {Juan Julián Merelo Guervós}\\[2cm]
\includegraphics[width=0.3\textwidth]{imagenes/etsiit_logo.png}\\[0.1cm]
\textsc{Escuela Técnica Superior de Ingenierías Informática y de Telecomunicación}\\
\textsc{---}\\
Granada, algún mes de 2018
\end{minipage}
\end{titlepage}

	
	% Plantilla prefacio UGR
	\chapter*{}
\thispagestyle{empty}

\begin{center}
{\large\bfseries Chief \\ Apoyo a la labor de la Policía Local}\\
\end{center}
\begin{center}
Sergio Cervilla Ortega\\
\end{center}

%\vspace{0.7cm}

\noindent{\textbf{Resumen}}\\

Lorem ipsum dolor sit amet, consectetur adipisicing elit, sed do eiusmod tempor incididunt ut labore et dolore magna aliqua. Ut enim ad minim veniam, quis nostrud exercitation ullamco laboris nisi ut aliquip ex ea commodo consequat. Duis aute irure dolor in reprehenderit in voluptate velit esse cillum dolore eu fugiat nulla pariatur. Excepteur sint occaecat cupidatat non proident, sunt in culpa qui officia deserunt mollit anim id est laborum.

\cleardoublepage

\thispagestyle{empty}

\noindent\rule[-1ex]{\textwidth}{2pt}\\[4.5ex]

D. \textbf{Juan Julián Merelo Guervós}, Profesor del Área de XXXX del Departamento de Arquitectura y Tecnología de Computadores de la Universidad de Granada.

\vspace{0.5cm}

\textbf{Informo:}

\vspace{0.5cm}

Que el presente trabajo, titulado \textit{\textbf{Chief, blah blah}},
ha sido realizado bajo mi supervisión por \textbf{Sergio Cervilla Ortega}, y autorizo la defensa de dicho trabajo ante el tribunal
que corresponda.

\vspace{0.5cm}

Y para que conste, expiden y firman el presente informe en Granada a X de mes de 2018.

\vspace{1cm}

\textbf{El director: }

\vspace{5cm}

\noindent \textbf{Juan Julián Merelo Guervós}

	
	% Índice de contenidos
	\newpage
	\tableofcontents

	% Índice de imágenes y tablas
	\newpage
	\listoffigures
	\newpage

	% Introducción 
	\chapter{Introducción}

Vivimos cada día en un mundo más conectado a Internet, es una época en la que se está dejando de 
lado el lápiz y el papel. Donde se está abriendo paso la tecnología. Estos nuevos medios, bien
usados, son capaces de facilitarnos en numerosas ocasiones las tareas del día a día. Convierte acciones
tediosas en simples gestos que, con el tiempo, acabamos automatizando y realizamos sin esfuerzo. 
Bajo dicha motivación, decidí crear un sistema informatizado para ayudar en las labores del cuerpo policial de Maracena.\\

Previamente, el cuerpo policial utilizaba un esquema bastante anticuado para desemplear sus funciones.
Por ejemplo, los registros de incidencias se iban apuntando manualmente y al final del turno
se reunían para poner en común todos los datos en un documento que posteriormente se guardaba.
Y al igual que en el ejemplo anterior, debían esperar a llegar a las dependencias para poder 
rellenar cualquier tipo de denuncia administrativa con los datos que habían apuntado en 
un papel para no olvidarlos.\\

Para ello, he diseñado una aplicación web que se encarga de gestionar una gran cantidad de 
modelos de denuncia y que permite que se rellenen de forma segura, a prueba de 
fallos y de una manera muy intuitiva para el usuario. Además de poder gestionar dichos documentos,
se ha creado un sistema de registro de incidencias completamente automatizado para que en 
ningún momento se pierda ningún dato y un gran número de funcionalidades más que se explicarán en 
posteriores puntos. El Ayuntamiento de Maracena será pionero en este ámbito, ya que han sido los 
primeros que han apostado por las nuevas tecnologías. Una gran apuesta, pero que traerá grandes
beneficios a la calidad del trabajo del cuerpo policial y a la seguridad de los ciudadanos. \\

Este Trabajo de Fin de Grado ha sido realizado completamente por mi, Sergio 
Cervilla Ortega durante el año 2018 y es libre bajo la licencia GNU General Public License v3.0 \cite{gplv3}\\ 

Ha sido desarrollado bajo el permiso, aceptación, revisión y ayuda del Ayuntamiento de Maracena, localidad en
la que resido a fecha de hoy.\\



	% Descripción del problema y hasta donde se llega
	\chapter{Descripción del problema}
Los objetivos del desarrollo de este proyecto son los siguientes:
\begin{enumerate}
    \item \textbf{Facilitar el trabajo a los policías locales de Maracena.}\\
    Debido a la inclusión de un sistema informatizado para la elaboración de 
    ficheros administrativos, gestión de incidentes y un sistema de croquis. Consiguiendo,
    por tanto, una mejora en la productividad de los agentes y un aumento de la seguridad
    global de los datos almacenados.
    
    \item \textbf{Inclusión del software libre en organismos del estado.}\\
    Dejando a un lado herramientas privativas sobre las que no tenemos un pleno control
    de los datos que están analizando y además fomentando el desarrollo libre. Porque 
    de esta manera cualquier persona puede sumarse al desarrollo y mejora de este software.

    \item \textbf{Elaboración de un entorno de pruebas real.}\\
    Se persigue la implementación de un entorno virtual lo más parecido posible a la 
    realidad para que el uso de la aplicación pueda ser probado antes del despliegue 
    final. 
   
    \item \textbf{Probar los conocimientos adquiridos a lo largo del grado.}\\
    Creando un sistema completo en el que tendremos que tener en cuenta todos los 
    aspectos técnicos adquiridos en el transcurso de la carrera. Dada la embergadura del
    proyecto, es necesario disponer de una base muy consolidada de los conocimientos adquiridos
    previamente.

    \item \textbf{Aprendizaje de tecnologías punteras en el sector.}\\
    Durante el desarollo del proyecto se ha buscado aprender y utilizar tecnologías en 
    auge que tienen un alto potencial. Esta decisión se toma en base  a la cantidad de gente
    y empresas que las utilizan, la comunidad tan amplia que tienen así como las contínuas
    mejoras que están recibiendo.  

\end{enumerate}



	% Estado del arte
	% 	1. Crítica al estado del arte
	% 	2. Propuesta
	\chapter{Estado del arte}

Actualmente el estado del mercado para este tipo de desarrollos está muy limitado dado que 
la mayoría de los datos que manejan dichos sistemas son de carácter sensible. Por tanto, se 
debe tener un especial cuidado con el tratamiento de los mismos tanto a la hora de transmitirlos
como en su almacenamiento.\\

Esta premisa ha conseguido que la única aplicación actual del mercado, y aún en fase de desarrollo,
utilizada por la Guardia Civil de Tráfico sea de código privativo y una caja negra para el usuario.
De hecho, está en una fase tan temprana de su desarrollo que aún carece de nombre y de prácticamente 
utilidad.\\ 

insertar más datos de la GC

% [hablar de fuera de españa]

% [hablar del open source]

	% Análisis del problema
	% 1. Análisis de requisitos
	% 2. Análisis de las soluciones
	% 3. Solucion propuesta
	% 4. Análisis de seguridad

	% Diseño de la solución
	\chapter{Tecnologías}

\subsection{Node.js}

Node.js es un entorno de ejecución para JavaScript construido con el motor de JavaScript
V8 de Chrome. En este proyecto se ha utilizado para elaborar el servidor de la aplicación.\\

Node.js tiene una característica que lo diferencia del resto de posibilidades para implementar 
el servidor y es que utiliza un modelo de operaciones E/S sin bloqueo y orientado a eventos 
asícronos. Esto nos ofrece la la posibilidad para que cuando se detecte una conexión se atienda
pero el resto del tiempo estará inactivo. Con este nuevo modelo conseguimos servidores altamente 
escalables de una manera muy sencilla, rápida y eficiente.\\

Este diseño contrasta directamente con el modelo con el modelo de concurrencia más común hoy en 
día en el que se usan hilos del Sistema Operativo. Las operaciones en red basadas en este modelo
son relativamente ineficientes además de complejas en su uso. Cabe mencionar que, aunque Node.js 
no utilice hilos del Sistema Operativo implique que no podamos aprovechar los múltiples cores
de nuestro sistema.\\

Otro punto a favor de Node.js es que no tendremos que preocuparnos por la posibilidad de un 
bloqueo en el proceso debido a que no existe dicha posibilidad por la naturaleza de dicha 
tecnología, la cual está basada en un bucle de eventos.

\subsection{MongoDB}

MongoDB es un sistema de base de datos multiplataforma y orientado a documentos de esquema libre.
Esto quiere decir que cada entrada puede tener un esquema de datos diferente al anterior, por
lo que podemos tener entradas que difieran en el número de atributos entre sí.\\

Debido a esta característica los datos no son guardados en registros, como en las bases SQL, 
sino que se guardan en documentos. Estos documentos son del tipo BSON, que es una representación
binaria de los archivos JSON.\\

Esta tecnología ha sido escrita en C++, por lo que está en contacto con el \textit{bare metal}. Esto
último es realmente importante ya que le permite optimizar los recursos y acceder a ellos de una 
manera extremadamente eficiente, con lo que consigue una velocidad muy alta en todas sus operaciones.\\

Cabe mencionar que MongoDB es una plataforma completamente distribuida, lo que nos proporciona 
un nuevo nivel de escalabilidad y disponibilidad. Esto quiere decir que a medida que nuestro desarrollo
crezca tanto en volumen de datos como en rendimiento MongoDB se escala automáticamente sin tiempos
de caída y sin cambios en nuestra aplicación.\\

Por último, cabe mencionar que la licencia de esta tecnología es GNU AGPL 3.0, por lo que se trata de un software de licencia
libre.

\subsection{Vagrant}

Vagrant es una herramienta para construir y administrar entornos de máquinas virtuales de una 
manera sencilla. Debido a que está diseñado para ser sencillo de utilizar y a que se centra en 
la automatización hace que utilicemos el menor tiempo posible en la preparación del entorno virtual
que necesitamos para comenzar a trabajar.\\

Nos proporciona un entorno muy rápido y simple de configurar, reproducible y portable utilizando
tecnología puntera en el sector de la virtualización como puede ser VirtualBox, VMWare, AWS, Docker... 


\subsection{Ansible}





	% Desarrollo bajo sprints: 
	% 	1. Permitir registros y login de usuarios
	% 	3. Desarrollo del sistema de incidencias
	% 	2. Desarrollo del sistema de denuncias administrativas
	% 	3. Desarrollo del sistema de croquis

	% Resultados

	% Presupuesto

	% Conclusiones

	% Trabajos futuros


	
	\newpage
	\bibliography{bibliografia}
	\bibliographystyle{plain}
	
\end{document}

